\begin{figure}[H]
\centering
\begin{minipage}[t]{0.43\textwidth}
  \vspace{0pt}\centering
  \includegraphics[width=\linewidth]{figures/ch06/panels/fig6_4_scale_1.pdf}\\[-0.2em]
  {\small 大噪声：粗尺度几何}
\end{minipage}\hfill
\begin{minipage}[t]{0.43\textwidth}
  \vspace{0pt}\centering
  \includegraphics[width=\linewidth]{figures/ch06/panels/fig6_4_scale_2.pdf}\\[-0.2em]
  {\small 中间尺度}
\end{minipage}

\vspace{0.8em}

\begin{minipage}[t]{0.43\textwidth}
  \vspace{0pt}\centering
  \includegraphics[width=\linewidth]{figures/ch06/panels/fig6_4_scale_3.pdf}\\[-0.2em]
  {\small 小噪声：局部细节}
\end{minipage}\hfill
\begin{minipage}[t]{0.43\textwidth}
  \vspace{0pt}
  \caption{\textbf{不同噪声水平对应不同的得分场。}
  同一个双分量分布在三个尺度上用高斯噪声进行平滑。大噪声揭示粗尺度的全局几何；随着噪声减小，分离的峰和更尖锐的局部结构逐渐显现。因此，多尺度得分模型以噪声水平为条件，而不是让所有尺度共享同一个得分场。为便于阅读，箭头长度经过归一化。}
  \label{fig:score-multiscale-fields}
\end{minipage}
\end{figure}
